\section{Generalized measurements and decoherence (9/23)}

\referto{\nielsen~2.2.6, 2.4.2, 8.2; \tong~16.4; \anastopolous~8.2}

\subsection{Generalized measurements on a mixed state}

Consider the product state $\rho_s\otimes\rho_a$ as the initial state of a composite system $\mathbb{V}_s\otimes\mathbb{V}_a$, consisting of a physical system $s$ unentangled with the experimental apparatus $a$. We assume the apparatus density matrix $\rho_a$ has the following spectral decomposition:
\begin{align}
    \rho_a=\sum_jp_j\dket{j}_a\dbra{j}_a,~~~\sum_jp_j=1,~0<p_j\leq1.
\end{align}
where $\{\dket{j}_e\}$ is a complete set of orthonormal basis for the apparatus Hilbert space $\mathbb{V}_e$. 

The system and apparatus evolve together under a unitary operator $\hat U$, before a projective measurement is performed for an observable (i.e. Hermitian operator) on the apparatus Hilbert space
\begin{align}
    \hat O=\hat 1_s\otimes(\sum_qq~\dket{q}_a\dbra{q}_a)
\end{align}
It's straightforward to show that for measurement outcome $q$, the initial state $\rho_s$ of the physical system collapses into the following (unnormalized) state
\begin{align}
    \rho_q\propto\sum_j M_{q,j}\rho_sM^\dagger_{q,j}
\end{align}
with a probability distribution
\begin{align}
    \text{Pr}(q)=\sum_j\text{Tr}(M_{q,j}\rho_sM^\dagger_{q,j})
\end{align}
where the measurement operators $\{M_{q,j}\}$ are defined as
\begin{align}
    M_{q,j}\equiv\expval{q|\hat U|j}\sqrt{p_j}
\end{align}
satisfying the completeness relation
\begin{align}
    \sum_{q,j}M^\dagger_{q,j}M_{q,j}=\hat1_s
\end{align}
The above discussion is completely in parallel to section \ref{sec:generalized measurements}, merely generalizing everything to the context of mixed states. Therefore we can formulate the measurement postulate of quantum mechanics for density matrices:

\textbf{Postulate 2b}: A \emph{generalized quantum measurement} is described by a collection $\{M_m\}$ of measurement operators satisfying the completeness relation:
\begin{align}
    \sum_m M_m^\dagger M_m=\hat 1.
    \end{align}
If the state of the quantum system is $\rho$ immediately before the measurements, then the probability that measurement result $m$ occurs is given by
    \begin{align}
    \text{Pr}(m)=\text{Tr}(\rho M_m^\dagger M_m)
    \end{align}
    and the (unnormalized) state of the system after the measurement is $\rho_m\propto M_m\rho M_m^\dagger$. 

In a generalized measurement, the set of positive operators $\{E_m\equiv M_m^\dagger M_m\}$ is known as a POVM (positive operator-value measure). A generalized measurement is sometimes also known as a POVM measurement. 

\subsection{Examples of generalized measurements}

(1) \emph{Quantum Zeno effect}

In a single-qubit system, consider a repeated 2-step procedure where the 1st step of unitary time evolution
for a duration $T/N$ is followed by the 2nd step of a projective measurement on $\hat X$:
\begin{align}
    \hat X=\hat P_+-\hat P_-,~~~\hat P_+=\proj{+},~\hat P_-=\proj{-}.
\end{align}
where
\begin{align}
    \dket{\pm}=\frac1{\sqrt{2}}(\dket{0}\pm\dket1)
\end{align}
are the eigenstates of Pauli $\hat X$ operator. 

In this example, the measurement operators associated with the $\pm$ outcome after
each step are $\{\hat M_\pm\}$ defined as follows:
\begin{align}
    \hat M_\pm=\hat P_\pm\,U(T/N),~~~M^\dagger_+M_++M^\dagger_-M_-=\hat1.
\end{align}


(2) \emph{Detection of a photon}

Let $\dket{0}$ denote no photon and $\dket{1}$ denote one photon.  A simple
photon detection model uses
\begin{align}
    \hat M_0=\proj{0},~~~\hat M_c=\dket{0}\dbra{1}
\end{align}
where the subscript $c$ denotes a click in the photon reader. The photon is annihilated after the click outcome. It's straightforward to check the completeness relation holds for $\{M_0,M_c\}$. 

(3) \emph{Distinguishing quantum states}

Suppose a qubit is prepared in either
\[
  \dket{0}
  \qquad\text{or}\qquad
  \dket{+}=\frac{1}{\sqrt{2}}\bigl(\dket{0}+\dket{1}\bigr).
\]
The task is to determine which state was received.

(i) {Projective measurement in the $Z$ basis}:

The measurement operators are projection operators
\[
  \hat P_0=\proj{0},
  \qquad
  \hat P_1=\proj{1}.
\]
If the outcome is $1$, the received state must have been $\dket{+}$ because
$\expval{1|0}=0$.  If the outcome is $0$, the two possibilities cannot be
distinguished.

(ii) {Projective measurement in the $X$ basis}:

The measurement operators are again projectors
\[
  \hat P_+=\proj{+},
  \qquad
  \hat P_-=\proj{-}.
\]
If the outcome is $+$, the states cannot be distinguished.  If the outcome
is $-$, the received state must have been $\dket{0}$ because
$\expval{-|+}=0$.

(iii) {A three-outcome POVM}:

Consider the POVM $\{E_m=M_m^\dagger M_m\}$:
\begin{align*}
  \hat E_1&=\frac{1}{2}\proj{1}
  =\begin{pmatrix}0&0\\[2pt]0&\tfrac12\end{pmatrix},\\
  \hat E_2&=\frac{1}{2}\proj{-}
  =\frac14\begin{pmatrix}1&-1\\[2pt]-1&1\end{pmatrix},\\
  \hat E_3&=\hat1-\hat E_1-\hat E_2
  =\begin{pmatrix}\tfrac34&\tfrac14\\[2pt]\tfrac14&\tfrac14\end{pmatrix}
  \ge 0.
\end{align*}
The outcomes have the following interpretation:
\begin{enumerate}[label=\arabic*.]
  \item Outcome $1$: the state must have been $\dket{+}$.
  \item Outcome $2$: the state must have been $\dket{0}$.
  \item Outcome $3$: the measurement is inconclusive.
\end{enumerate}




\subsection{Measurements induce decoherence}


In its eigenbasis $\{\dket{j}|1\leq j\leq\text{dim}\mathbb{V}\}$, a mixed state (\ref{ensemble of state}) is an ensemble of pure states $\{\dket{j}\}$
\begin{align}\label{mixed state}
\hat\rho=\sum_{j}p_j\dket{j}\dbra{j},~~~\sum_jp_j=1.
\end{align}
while a generic pure state in the same physical system can be written as a (normalized) superposition of the basis vectors $\{\dket{j}\}$:
\begin{align}\label{pure state}
\dket{\psi}=\sum_j\psi_j\dket{j},~~~\sum_j|\psi_j|^2=1.
\end{align}
But there is a big difference between the (``incoherent'') ensemble summation in the mixed state (\ref{mixed state}) and the (``coherent'') superposition in the pure state (\ref{pure state}), in the sense that the pure state contains more information in the relative phases of coefficients $\psi_j$ that can be measured experimentally. To make it clear, we use the normalization condition in (\ref{pure state}) to define the probability distribution $p_j\equiv|\psi_j|^2$, and the normalized pure state can be rewritten as
\begin{align}\label{pure state:normalized}
\dket{\psi}=\sum_j\sqrt{p_j}e^{\imth\phi_j}\dket{j},~~~\sum_jp_j=1.
\end{align}
While the global phase factor is a redundancy with no physical significance, there are $(\text{dim}\mathbb{V}-1)$ relative phases that can be detected by physical observables. 

To illustrate, consider a single qubit example, where we use $(\hat X,\hat Y,\hat Z)$ to label the 3 Pauli matrices. We consider a mixed state $\hat\rho$ and a pure state $\psi$, both having $1/2$ probability for either $Z=+1$ or $Z=-1$. In the computational basis $\dket{0}\equiv\dket{\uparrow},~\dket{1}\equiv\dket{\downarrow}$ that diagonalizes $\hat Z$, the mixed state is written as
\begin{align}
\hat\rho=\frac12(\dket{0}\dbra{0}+\dket{1}\dbra{1})=\frac12\bpm1&0\\0&1\epm
\end{align}
and the pure state as
\begin{align}
\dket{\psi}=\frac1{\sqrt2}(\dket{0}+e^{\imth\phi}\dket{1})=\frac1{\sqrt2}\bpm1\\e^{\imth\phi}\epm
\end{align}
If we now measure another observable $\hat X$ for both states, we end up with different expectation values:
\begin{align}
\trace(\hat\rho\hat X)=0,~~~\expval{\psi|\hat X|\psi}=\cos\phi.
\end{align}
exactly rooted in the relative phases and associated quantum inteferences in the pure state.

As a result, the reduction from a pure state (\ref{pure state:normalized}) to a mixed state (\ref{mixed state}) is a process of decoherence, i.e. a reduction of quantum interference. This can be achieved by a projective measurement (see postulate \ref{postulate 3:projective measurements}), where the measurement outcome is unknown. Consider the implementation of a measurement for physical observable diagonal in basis $\{\dket{j}\}$
\begin{align}
\hat O\equiv\sum_j O_j\dket{j}\dbra{j}
\end{align}
on the pure state (\ref{pure state:normalized}), If we don't know the outcome, after a projective measurement, we should view the resulting state as a summation of projector $\dket{j}\dbra{j}$ weighted by probability $p_j=|\expval{j|\psi}|^2$, and hence we end up with the mixed state (\ref{mixed state}). This means a quantum measurement will generally decohere a quantum state, even if we don't know the measurement outcome. This is in sharp contrast to a classical system. 


